% ---------------------------------------------------------------------
%  APÊNDICE C — COMO REPRODUZIR OS RESULTADOS
%  Revisão de 3/11, item 5.2: o caminho mínimo do clone às tabelas e
%  figuras. O README do repositório tem a instalação detalhada.
% ---------------------------------------------------------------------
\chapter{Como reproduzir os resultados}
\label{ap:reproduzir}

Todo o código do trabalho está no repositório
\url{https://github.com/BonumLusor/ardupilot-custom-control}, sob a licença
GPLv3, a mesma do ArduPilot. Isso inclui os backends dos controladores, o
injetor de distúrbios, os cenários, a modelagem, as análises e este texto. Os
resultados dos Capítulos~\ref{cap:modelagem} a~\ref{cap:resultados} saíram da
versão \pendente{etiqueta e DOI do Zenodo da versão da defesa}. Este apêndice
dá o caminho mínimo, do clone às tabelas e figuras. O \arquivo{README.md} do
repositório traz os detalhes da instalação e os problemas conhecidos.

\section{Ambiente}

O trabalho rodou no Ubuntu 24.04, sob o WSL2, com Python 3.12. O ArduPilot
entra como submódulo, fixo no \textit{commit} \arquivo{ce9c2d01a4}
(4.8-dev). Os backends LQR e LQG ficam em \arquivo{firmware/}, fora do
submódulo. O \arquivo{firmware/instalar.sh} os liga ao submódulo por
link simbólico e aplica o \arquivo{firmware/ardupilot.patch}. O patch traz as
edições do \textit{frontend}, a mensagem de log \arquivo{CCTQ} com a atitude
verdadeira e a contagem do \arquivo{SIM_TWIST_TIME} em microssegundos.

\begin{lstlisting}[language=bash,numbers=none]
git clone --recurse-submodules \
    https://github.com/BonumLusor/ardupilot-custom-control.git TCC
cd TCC
ardupilot/Tools/environment_install/install-prereqs-ubuntu.sh -y
python3 -m venv ~/venv-ardupilot
~/venv-ardupilot/bin/pip install -r requirements.txt
firmware/instalar.sh
cd ardupilot && ./waf configure --board sitl && ./waf copter && cd ..
\end{lstlisting}

O \arquivo{run_ensaio.py} chama o \arquivo{sim_vehicle.py} com
\arquivo{--no-rebuild}, e todos os voos usam o binário compilado antes. Por isso,
depois de qualquer mudança no firmware, é preciso compilar de novo com
\codigo{./waf copter}. Cada voo grava em \arquivo{versao.json} o hash do binário
do SITL, o do injetor e os \textit{commits}. O lote da monografia usou o binário
cujo hash está em \arquivo{ensaios/lotes/binario_lote.sha256}.

\section{Um ensaio}

Um ensaio é um cenário (\arquivo{ensaios/cenarios/}) voado com um controlador
(\arquivo{ensaios/controladores/}):

\begin{lstlisting}[language=bash,numbers=none]
source ~/venv-ardupilot/bin/activate
ensaios/run_ensaio.py --cenario rolagem/torque_impulso --controlador cc_lqg
\end{lstlisting}

O resultado fica em
\arquivo{ensaios/resultados/rolagem/torque_impulso/cc_lqg/<carimbo>/}. Lá
estão o log \arquivo{.BIN}, os parâmetros e as métricas na atitude do AHRS
(\arquivo{metricas.json}). Há também a série reamostrada a \SI{400}{\hertz}
(\arquivo{series.csv}), o painel do voo e o \arquivo{versao.json}. As métricas na
atitude verdadeira (\arquivo{metricas_verdade.json} e
\arquivo{serie_verdade.csv}) saem de \codigo{python -m analise.verdade}, rodado
em \arquivo{ensaios/}. Com \codigo{-{}-repeticoes N} ou \codigo{-{}-sementes 1-5}, o
mesmo par voa várias vezes. Com \arquivo{--pular-feitos}, o que já foi voado com
o mesmo binário é pulado.

\section{O lote, as tabelas e as figuras}

O lote da monografia está em \arquivo{ensaios/lotes/lote_revisao.sh}. São
cerca de 560 voos, em umas oito horas. Ele começa pelo teste de determinismo e
pode ser retomado: rodá-lo de novo pula os voos feitos. Depois dele, as
métricas na atitude verdadeira, as tabelas e as figuras dos
Capítulos~\ref{cap:modelagem} a~\ref{cap:resultados}:

\begin{lstlisting}[language=bash,numbers=none]
ensaios/lotes/lote_revisao.sh --l4-inerente
cd ensaios
python -m analise.verdade resultados --binario $(cat lotes/binario_lote.sha256)
./tabelas_monografia.py
./tabelas_monografia.py --atitude ahrs --destino /tmp/tabelas_ahrs
./figuras_monografia.py
\end{lstlisting}

As tabelas e figuras usam só os voos do binário do lote. A base de atitude
padrão é a verdadeira, e cada tabela grava, ao lado, a nota com a base e o
número de voos de cada célula.

\section{Modelo, projeto e confronto com o SITL}

O modelo do SITL, o projeto dos controladores e o confronto entre modelo e SITL
(Capítulos~\ref{cap:modelagem} e~\ref{cap:projeto} e
Apêndice~\ref{ap:validacao-proposto}) saem de \arquivo{modelagem/}:

\begin{lstlisting}[language=bash,numbers=none]
cd modelagem
./modelo_sitl.py --parm <params.parm de um voo>
./projeto.py --figuras ../monografia/figuras --tabelas ../monografia/tabelas
./validacao.py --binario $(cat ../ensaios/lotes/binario_lote.sha256) \
    --figuras ../monografia/figuras --tabela ../monografia/tabelas/validacao.tex
\end{lstlisting}

O \arquivo{projeto.py} refaz a varredura e o portão da
Seção~\ref{sec:projeto-lqg}; com \arquivo{--cache}, reaproveita a varredura já
calculada. Os ganhos que ele escreve nos YAMLs de \arquivo{ensaios/controladores/}
são os que o lote voa.
